% Part: proof-theory
% Chapter: sequent-calculus
% Section: rules-proofs

\documentclass[../../../include/open-logic-section]{subfiles}

\begin{document}

\olfileid[pt]{pt}{seq}{rp}

\olsection{Regras e \usetoken{P}{proof}}

Começamos apresentando algumas definições e fixando certa terminologia.

Considere, por exemplo, as duas regras de \Log{G3c} que permitem
!!{proving} sequentes que contêm $!A \land !B$ a partir de sequentes que
contêm $!A$ e~$!B$:
\[
\Axiom$ !A, !B, \Gamma \fCenter \Delta$
\RightLabel{\LeftR{\land}}
\UnaryInf$ !A \land !B, \Gamma \fCenter \Delta$
\DisplayProof
\qquad
\Axiom$\Gamma \fCenter \Delta, !A$
\Axiom$ \Gamma \fCenter \Delta, !B$
\RightLabel{\RightR{\land}}
\BinaryInf$ \Gamma \fCenter \Delta, !A \land !B $
\DisplayProof
\]
Os sequentes acima da linha são chamados de \emph{premissas}, e o sequente
abaixo da linha, de \emph{conclusão}. As !!{formula}s em $\Gamma$ e
$\Delta$ em cada regra são chamadas de \emph{contexto} ou de
\emph{!!{formula}s laterais}. A !!{formula} na conclusão que não faz parte
do contexto é chamada de !!{formula} \emph{principal}, e as !!{formula}s
nas premissas que não pertencem ao contexto são chamadas de !!{formula}s
\emph{ativas} (ou \emph{auxiliares}).

As regras dos quantificadores são um pouco mais complicadas. P.ex., as
regras de \Log{G1c} para $\lforall$ são:
\[\Axiom$ !A(t), \Gamma \fCenter \Delta$
\RightLabel{\LeftR{\lforall}}
\UnaryInf$ \lforall[x][!A(x)],\Gamma \fCenter \Delta$
\DisplayProof
\qquad
\Axiom$ \Gamma \fCenter \Delta, !A(a) $
\RightLabel{\RightR{\lforall}}
\UnaryInf$ \Gamma \fCenter \Delta, \lforall[x][!A(x)]$
\DisplayProof
\]
Na regra \LeftR{\lforall}, a !!{formula} ativa é~$!A(t)$, onde $t$ é um
termo fechado, O.s.\@, um termo sem variáveis. Tal inferência é correta---ou
seja, corresponde à regra esquemática \LeftR{\lforall}---se existem alguma
!!{formula}~$!A$, alguma variável~$x$ e algum termo fechado~$t$ tais que a
!!{formula} principal na conclusão é $\lforall[x][!A]$ e a !!{formula}
ativa na premissa é~$\Subst{!A}{t}{x}$. A regra \RightR{\lforall} funciona
do mesmo modo, exceto que, em vez de um termo~$t$ qualquer, a !!{formula}
ativa na premissa deve ser~$\Subst{!A}{a}{x}$, onde $a$ é !!a{constant}.
Essa constante é chamada de \emph{variável própria} da inferência
\RightR{\lforall}. A inferência só é correta se o sequente inferior não
contiver~$a$, isto é, se $a$ não ocorrer em $\Gamma$, $\Delta$ ou~$!A$.
Isso se chama \emph{condição da variável própria}.

Sistemas de sequentes como \Log{G1c} e \Log{G3c} consistem em uma coleção
dessas regras especificadas esquematicamente, juntamente com alguns
sequentes (também especificados esquematicamente) que contam como axiomas.
Os !!^{proof}s em tal sistema são árvores finitas de sequentes geradas
indutivamente a partir de axiomas pelo uso das regras de inferência. Cada
folha é um axioma, e todo outro sequente decorre dos sequentes imediatamente
acima dele por uma das regras de inferência do sistema.

\begin{defn}\ollabel{defn:proof}
Os \emph{!!^{proof}s} em um sistema de sequentes são árvores finitas de
sequentes definidas indutivamente da seguinte maneira:
\begin{enumerate}
  \item\ollabel{defn:prf:axiom} Todo axioma é !!a{proof}.
  \item\ollabel{defn:prf:one-prem} Se $\pi_1$ é !!a{proof} que termina em
  um sequente~$S_1$, e $S$ é um sequente que decorre de~$S_1$ por uma
  regra~$R$ de uma premissa, então
  \[
  \AxiomC{}
  \RightLabel{$\pi_1$}
  \DeduceC{$S_1$}
  \RightLabel{$R$}
  \UnaryInfC{$S$}
  \DisplayProof
  \]
  é !!a{proof}.
  \item\ollabel{defn:prf:two-prem} Se $\pi_1$ e $\pi_2$ são !!{proof}s que terminam nos sequentes~$S_1$ e~$S_2$, e $S$ é um
  sequente que decorre de~$S_1$ e$S_2$ por uma regra~$R$ de duas premissas, então
  \[
  \AxiomC{}
  \RightLabel{$\pi_1$}
  \DeduceC{$S_1$}
  \AxiomC{}
  \RightLabel{$\pi_2$}
  \DeduceC{$S_2$}
  \RightLabel{$R$}
  \BinaryInfC{$S$}
  \DisplayProof
  \]
  é !!a{proof}.
\end{enumerate}
\end{defn}

\Olref{defn:proof} define !!{proof}s como árvores de sequentes. Imaginamos
que essas árvores crescem ``para cima''. Os nós superiores (folhas) dessa
árvore são axiomas e são chamados de \emph{sequentes iniciais}. O sequente
inferior é chamado de \emph{sequente final}. Se $\pi$ é !!a{proof} com
sequente final~$S$, também escrevemos $\pi \Proves S$. Se existe
!!a{proof} de um sequente~$S$, escrevemos $\Proves S$, podendo indicar à
esquerda do símbolo~$\Proves$ o sistema no qual a demonstração existe; p.ex., $\Log{G1c} \Proves !A, !B \Sequent !A \land !B$. Todo sequente em
!!a{proof}, com exceção dos sequentes iniciais, é a \emph{conclusão} de uma
inferência por alguma regra~$R$. O único ou os dois sequentes imediatamente
acima de um sequente em !!a{proof} (O.s.\@, os pais do nó) são chamados de
\emph{premissas} da inferência.

Os !!{proof}s usados na construção indutiva de !!a{proof} são chamados de
seus \emph{sub-!!{proof}s}. Os sub-!!{proof}s que terminam nas premissas de
uma inferência são chamados de \emph{sub-!!{proof}s imediatos} da
inferência.

Muitas vezes é necessário medir a complexidade dos !!{proof}s por um número
natural. Poderíamos simplesmente contar o número de sequentes em !!a{proof},
mas uma medida frequentemente mais útil é a !!{height} de !!a{proof}: ela é
o número máximo de sequentes entre o sequente final e um sequente inicial.
Também podemos dar uma definição indutiva.

\begin{defn}
A \emph{!!{height}}~$\pheight{\pi}$ de !!a{proof}~$\pi$ é definida
indutivamente da seguinte maneira.
\begin{enumerate}
  \item Se $\pi$ consiste em um único axioma, $\pheight{\pi} = 0$.
  \item Se $\pi$ termina em uma inferência com uma premissa cujo
  sub-!!{proof} imediato é~$\pi_1$, então $\pheight{\pi} = \pheight{\pi_1} + 1$.
  \item Se $\pi$ termina em uma inferência com duas premissas cujos
  sub-!!{proof}s imediatos são~$\pi_1$ e~$\pi_2$, então $\pheight{\pi} =
  \max(\pheight{\pi_1}, \pheight{\pi_2}) + 1$.
\end{enumerate}
Se um sequente~$S$ tem !!a{proof} de altura $\le n$, escrevemos $\Proves[n] S$.
\end{defn}

\begin{ex}
Em \Log{G3c}, todo sequente da forma $!C, \Gamma \Sequent \Delta, !C$ é um
axioma, onde $!C$ deve ser atômica. Suponha que $!D$ e $!E$ sejam sentenças
atômicas. Então, $!D, !E \Sequent !D$ e $!D, !E \Sequent !E$ são instâncias
do axioma $!C, \Gamma \Sequent \Delta, !C$. P.ex., se tomarmos
$!C \ident !E$, $\Gamma = \{!D\}$ e $\Delta = \emptyset$ em
$!C, \Gamma \Sequent \Delta, !C$, obteremos exatamente o sequente
$!D, !E \Sequent !E$. (Lembre-se de que estamos lidando com
multi\emph{conjuntos}; portanto, $!D, !E \Sequent !E$ e
$!E, !D \Sequent !D$ são o mesmo sequente.)

Como $!D, !E \Sequent !D$ é uma instância de um axioma de~\Log{G3c}, ele
(por si só) conta como !!a{proof} em~\Log{G3c}; o mesmo vale para
$!D, !E \Sequent !E$ por \olref{defn:proof}\olref{defn:prf:axiom}.
Chamemos essas demonstrações de $\pi_1$ e~$\pi_2$. Por
\olref{defn:prf:two-prem}, podemos tomar esses dois !!{proof}s e gerar um
novo ($\pi_3$) usando a regra~$\RightR{\land}$ de duas premissas:
\[
\Axiom$!D, !E \fCenter !D$
\Axiom$!D, !E \fCenter !E$
\RightLabel{\RightR{\land}}
\BinaryInf$!D, !E  \fCenter !D \land !E $
\DisplayProof
\]
De $\pi_3$, por sua vez, obtemos o !!{proof}~$\pi_4$
\[
\Axiom$!D, !E \fCenter !D$
\Axiom$!D, !E \fCenter !E$
\RightLabel{\RightR{\land}}
\insertBetweenHyps{\hspace{5ex}}
\BinaryInf$!D, !E  \fCenter !D \land !E $
\LeftSubproofLabel{$\pi_3$}
\RightLabel{\LeftR{\land}}
\UnaryInf$!D \land !E  \fCenter !D \land !E$
\RightSubproofLabel{$\pi_4$}
\DisplayProof
\]
usando a regra~$\LeftR{\land}$ de uma premissa (por
\olref{defn:proof}\olref{defn:prf:one-prem}). O sub-!!{proof} imediato da
última inferência em~$\pi_4$ é~$\pi_3$. Os sub-!!{proof}s imediatos da
inferência $\RightR{\land}$ são $\pi_1$ e~$\pi_2$. Os sub-!!{proof}s de
$\pi_4$ são $\pi_1$, $\pi_2$, $\pi_3$ e o próprio $\pi_4$. Temos
$\pheight{\pi_1} = \pheight{\pi_2} = 0$, $\pheight{\pi_3} = 1$ e
$\pheight{\pi_4} = 2$. Mostramos que
$\Log{G3c} \Proves[4] !D \land !E  \fCenter !D \land !E$ para quaisquer
$!D$ e~$!E$ atômicas.
\end{ex}

\end{document}
